\documentclass[masters,nl,book,11pt,redheadings]{OU}

\addbibresource{bib.bib} % here you include the bibliography file. 

\usepackage{lipsum}

\title{A very long and cumbersome title of the Master's thesis}
\shorttitle{Short title of the Master's thesis}
\author{Student Name}
\thesisyear{2026}
\credits{30}
\faculty{Faculty of Science}
\studyprogramme{Master of Science in Environmental Sciences}
\supervisor{Supervisor Name}

\abstract{%
This short text demonstrates the abstract page generated by the class.
Replace it with a concise account of the problem, method, principal results, and main conclusion of the thesis.
}

\samenvatting{%
Deze korte tekst demonstreert de Nederlandstalige samenvatting. Vervang deze tekst door een beknopte beschrijving van het probleem, de methode, de belangrijkste resultaten en de conclusie van de scriptie.
}

\acknowledgements{%
This optional page can be used to acknowledge scientific, technical, and personal support received during the project.
}

\begin{document}

\startthesis

\chapter{Introduction}

The introduction establishes the scientific context, identifies the knowledge gap, and narrows the discussion towards a precise research question.
It should also explain why answering that question matters.

\lipsum[1-12]


\section{Research question}

State a question that can be answered with the data and methods available in the project.
Define the scope and the principal assumptions explicitly.

\chapter{Background}

Introduce the topic in a way that is accessible to a reader who is not an expert in the field.
Use citations to support statements and to indicate the origin of ideas.
This is also the place to introduce the theoretical framework and to define key concepts.

\lipsum[13-28]

\chapter{Literature review}

A literature review is a comprehensive survey of published sources on a specific topic.
It provides an overview of existing research, theories, and methodologies relevant to your study.
The purpose is to contextualise your research within the current state of knowledge, identify gaps in the literature, and establish how your work contributes to the field~\cite{Beck2005timeseries}.
A well-structured literature review critically analyses and synthesises existing work, rather than merely summarising it, and demonstrates your understanding of the subject matter and research landscape.

\lipsum[29-44]

\chapter{Results}

Present observations before interpreting them.
Tables and figures should each communicate one main point and remain understandable from their captions.

\lipsum[45-60]

\chapter{Discussion}

Relate the results to the research question and relevant literature.
Discuss uncertainty, limitations, and the extent to which the conclusions can be generalised.

\lipsum[61-76]

\chapter{Conclusion}

Answer the research question directly and briefly.
Do not introduce new data or new lines of argument in the conclusion.

\lipsum[77-90]

\references

\end{document}
